\section{Data}
\label{sec:data}

The analysis uses the harmonised Households Below Average Income (HBAI) extract of the Family Resources Survey, a repeated cross-section rather than a panel, covering 2016--2024, in 2023/24 prices, and excluding 2020/2021 due to COVID-19 disruption. The sample is restricted to children aged 16 or under living in England (the control group) or Scotland (the treatment group). Wales and Northern Ireland are excluded from the main design and used instead as falsification checks (Section~\ref{sec:results}). England contributes the large majority of observations for each year, while Scotland's smaller population share is reflected in correspondingly smaller cell counts, ranging from 537 (FYE2024) to 1,190 (FYE2017) children (Appendix Table~\ref{tab:summary_sample_composition}).

Following \citet{andersen2025} and \citet{nesom2025}, the headline design treats the November 2022 expansion as the single treatment date, with the post-treatment indicator refined to the exact 14 November cutoff within FYE2023 where household interview dates are available from the underlying FRS file; the earlier, age-limited 2021 rollout is not used to define treatment timing in the headline specification, since only a small under-6 subgroup is affected, limiting power to detect a separate age-cohort effect.

The primary outcome is the official DWP child material deprivation (MDCH) flag, a binary indicator equal to one if a child's household is classified as materially deprived under DWP's scoring methodology. For the pre-FYE2024 design used across most of the sample period, a household is asked whether it can afford each of 21 items (12 child-specific, nine whole-family); each item it cannot afford scores one, weighted by the item's prevalence in the wider population, and the household is classified as materially deprived once the summed, rescaled (0--100) score reaches 25 or more. 

The questionnaire was redesigned at FYE2024, replacing this 21-item, prevalence-weighted design with a revised 22-item (11 child, 11 household) simple-count design, in which every item carries equal weight and the threshold becomes four or more items lacked. Because this redesign breaks direct comparability with earlier years, this paper instead falls back on ten of the pre-FYE2024 design's twelve child-specific items (the two dropped for sparse or uncertain coverage), each recoded to 0/1 (deprived/not deprived); households surveyed under the redesigned FYE2024 questionnaire, which does not populate these ten items in the same form, are excluded from the resulting item-level measures rather than recoded onto a common scale. These ten items are also the basis for two additional composite measures, \emph{any deprivation} (at least one item lacking) and \emph{severe deprivation} (more than five items lacking, this paper's own construction rather than an official DWP statistic). No fixed-count threshold is well established in the literature for defining severity, and \citet{nicoriciuelliot2025} argue against arbitrary count thresholds altogether in favour of latent-class typologies (Section~\ref{sec:lit-measurement}), so these two composites are reported only as a robustness check in Appendix~\ref{sec:appendix-composites} rather than as headline outcomes. A secondary food-insecurity outcome, available only from FYE2020, is used for comparability with \citet{andersen2025} and as an additional pre-trend diagnostic. The deprivation module is administered on an identical basis across Great Britain, with no country-specific question wording or routing in the FRS instrument, so the measure is treated as consistently defined across the Scotland-England comparison.

Table~\ref{tab:summary_outcomes} summarises the four headline outcomes by country and period. The official MDCH flag rises in each country but markedly more in England (18.9\% to 26.5\%) than in Scotland (14.8\% to 15.9\%). Any deprivation -- lacking at least one of the ten items -- shows a sharper divergence still: it rises in England (36.4\% to 39.2\%) but falls in Scotland (33.4\% to 27.0\%), the clearest raw evidence motivating the difference-in-differences design below. Food insecurity, by contrast, rises by a similar margin in England but stays essentially flat in Scotland. Any and severe deprivation are examined further as a robustness check in Appendix~\ref{sec:appendix-composites}; this paper's main causal analysis focuses on the official MDCH flag and its item-level decomposition (Section~\ref{sec:methodology}).

\input{Tables/summary_outcomes.tex}

Table~\ref{tab:summary_background} reports weighted background characteristics for England and Scotland pooled across all years, replicating Table 2 of \citet{andersen2025}. Scotland has a smaller share of large (3+ child) families, a similar share of lone-parent households, a slightly younger and more female-headed household-head profile, and a substantially larger White household share (91.8\% vs.\ 76.6\%) than England. 

\input{Tables/summary_background.tex}

Table~\ref{tab:summary_a1_by_period} (Appendix~\ref{sec:appendix-a1-period}) decomposes these same characteristics by pre- and post-SCP period. Composition is broadly stable within each country, consistent with the identifying assumption that observable composition is not shifting sharply around the treatment date in a way that would itself explain a change in measured deprivation.

Two covariate sets are used throughout. A lean, six-variable set replicates the CASE evaluation's specification exactly \citep{andersen2025} (Table~\ref{tab:data_dictionary_lean}), deliberately excluding household income, benefit-receipt flags, and housing tenure, since none is part of CASE's published specification and income responds mechanically to receiving SCP while benefit receipt is itself an SCP eligibility criterion. A wide, approximately 35-variable set -- continuous income, housing and benefit-unit amounts, family composition, and all nine benefit-receipt flags -- is used only in the double machine learning (DML) analysis (Section~\ref{sec:results}), to test sensitivity to covariate breadth once the estimator no longer requires a hand-specified linear functional form.
 
All estimates are weighted by the dependent-child survey weight (\texttt{GS\_INDCH}). Standard errors are clustered at household level throughout, as discussed in Section~\ref{sec:lit-identification}. The analytic sample is approximately 47,000 child-year observations for the item-based outcomes and 51,000 for the official MDCH flag.
 
\section{Methodology}
\label{sec:methodology}
 
The empirical analysis proceeds in five completed stages: (1) baseline covariate-adjusted OLS difference-in-differences (DiD) on the single official outcome replicating \citeauthor{andersen2025}'s results; (2) the same design re-estimated jointly across all ten deprivation items in one household-clustered model; (3) a doubly robust machine-learning re-estimation of the baseline, testing sensitivity to functional form and covariate breadth; (4) the same machine-learning approach applied separately to each of the ten items; and (5) a jointly estimated, vector-outcome version of stage (4).
 
\subsection{Identification: assumptions and treatment allocation}
\label{sec:identification-assumptions}
 
The DiD estimates in this paper are interpreted as the causal intention-to-treat (ITT) effect of SCP on all children in Scotland, not a strict average treatment effect on the treated (ATT), since FRS/HBAI data cannot reliably identify SCP receipt and restricting to self-reported recipients would both shrink Scotland's already-small child sample (Appendix Table~\ref{tab:summary_sample_composition}) to an unworkably small size and introduce selection bias, as receipt correlates with the same low income that predicts higher deprivation. The resulting ITT estimate dilutes the true effect on genuine recipients roughly in proportion to the share of Scottish children outside SCP-eligible households, so the DiD coefficients reported in Section~\ref{sec:results} should be read as a lower bound on SCP's effect on the children it actually reaches, not the average effect on a randomly chosen Scottish child.

Interpreting this ITT effect causally requires, beyond parallel trends (Section~\ref{sec:parallel-trends}), correct treatment allocation, no anticipation, and no spillovers between Scotland and England. Scotland is the only UK nation where SCP is paid, making it the only possible treated unit, with England as the natural control under a shared, reserved welfare architecture across Great Britain. Wales and Northern Ireland, which never received SCP, are used instead as falsification checks (Section~\ref{sec:res-placebo}). Scotland also runs other Scotland-only, child-targeted transfers, notably the Best Start Grant and Best Start Foods (Section~\ref{sec:introduction}), so the Scotland$\times$Post coefficient more precisely captures Scotland's wider child-income-support reform package, of which SCP is the largest and only sharply-dated component. This is a limitation on causal interpretation partially addressed by the age-restricted robustness check in Section~\ref{sec:res-robustness}.

No anticipation requires that households did not adjust behaviour before SCP's rollout in expectation of receiving it, partially tested by the FYE2022 exclusion check (Section~\ref{sec:res-robustness}). No spillovers requires Scotland's receipt of SCP does not change England's outcomes. Cross-border migration is unlikely at these sample sizes, and Table~\ref{tab:summary_a1_by_period} shows no unusual shift in Scotland's household composition around the treatment date. Policy or media spillover is harder to rule out, since sustained UK-wide coverage of SCP could plausibly shift England's own counterfactual trajectory (Section~\ref{sec:introduction}); this is carried forward as an explicit limitation in the discussion.

\subsection{Identification: parallel trends}
\label{sec:parallel-trends}

Parallel pre-treatment trends are tested via Scotland$\times$time interactions, replicating the CASE evaluation's own diagnostic tables which include a covariate balance table and an outcome event-study table, jointly tested via a Benjamini--Hochberg-corrected Wald test. Standard errors are clustered at household level.

At annual level, and under Benjamini--Hochberg correction, every outcome and covariate in both diagnostic tables reads as consistent with parallel trends. Figure~\ref{fig:pt_composite} plots the underlying event-study coefficients. Pre-treatment estimates are statistically indistinguishable from zero, visually confirming the Wald test before a treatment effect opens up after the November 2022 rollout. The same pattern holds for food insecurity, reported in full in Appendix~\ref{sec:appendix-quarterly} (Figure~\ref{fig:pt_composite_grid}); the two composite outcomes examined as a robustness check (Appendix~\ref{sec:appendix-composites}) are likewise consistent with parallel trends at annual level, reported in the same figure.
 
\begin{figure}[htbp]
\centering
\includegraphics[width=0.7\textwidth]{Figures/es_MDCH_flag.png}
\caption[Parallel-trends event study, official MDCH flag (annual)]{Scotland$\times$year event-study coefficients, official MDCH flag, annual granularity, relative to the FYE2022 reference year. Blue is the pre-treatment period, red is post-treatment, with 95\% confidence bands; the dashed line marks the November 2022 rollout.}
\label{fig:pt_composite}
\end{figure}
 
A separate, more detailed check replicates the CASE paper's own event-study specification at quarterly level, using rolling 13-week windows anchored on the SCP rollout date rather than annual survey-year dummies. Table~\ref{tab:parallel_trends_quarterly_wald} (Appendix~\ref{sec:appendix-quarterly}) reports the resulting joint Wald tests, Benjamini--Hochberg corrected across all 14 outcomes and items tested. At this finer granularity, two individual items remain flagged even after correction, Bed/bedroom and Outdoor play area (Section~\ref{sec:disc-reconciliation}); the headline official MDCH flag remains clean at quarterly level. One of the two composite outcomes examined as a robustness check is also flagged at this finer granularity, discussed alongside its other results in Appendix~\ref{sec:appendix-composites}. The corresponding quarterly figure is reported in Appendix~\ref{sec:appendix-quarterly} (Figure~\ref{fig:pt_composite_quarterly}).
 
\subsection{Baseline difference-in-differences}
\label{sec:baseline-did}
 
The estimating equation for the single-outcome baseline is
\begin{equation}
Y_{it} = \alpha + \beta\,\text{Treated}_{i} + \delta\,(\text{Treated}_{i} \times \text{Post}_{t}) + X_{it}'\theta + \lambda_{t} + \varepsilon_{it},
\label{eq:baseline-did}
\end{equation}
fitted by weighted least squares with household-clustered standard errors, where $X_{it}$ denotes the covariate vector, which varies by specification as detailed below, $\lambda_{t}$ are survey-year fixed effects, and $\delta$, the coefficient on the Treated$\times$Post interaction, is the DiD estimate of SCP's intention-to-treat (ITT) effect, not a strict average treatment effect on the treated (ATT), for the reasons set out in Section~\ref{sec:identification-assumptions}. Three variants are reported: \emph{Simple} (no covariates beyond the DiD terms and year fixed effects), \emph{Adjusted} (the six CASE controls described in Section~\ref{sec:data}, the genuine \citeauthor{andersen2025} replication specification), and \emph{Extended} (Adjusted plus the child's own age, a labelled robustness variant, since several deprivation items are age-differentiated in what they capture). An item-by-item version of Equation~\ref{eq:baseline-did}, with \citet{benjaminihochberg1995} false-discovery-rate correction applied across the ten resulting tests, is also estimated, alongside an exact replication of \citet{andersen2025}'s own Table 3 specification (which includes an explicit \texttt{Post} main-effect term rather than absorbing it into year fixed effects, and reports the full covariate vector).

Equation~\ref{eq:baseline-did} omits a standalone \texttt{Post} main effect, relying on the year fixed effects $\lambda_{t}$ to absorb the common time trend instead. For most of the panel this is exact, since \texttt{Post} is a deterministic function of survey year; the one exception is FYE2023, within which the 14 November 2022 rollout cutoff falls, so households interviewed before and after that date share a year fixed effect despite differing in treatment timing. The CASE exact-replication specification (Table~\ref{tab:stage1_case_exact}) is identical to Adjusted in every other respect but adds \texttt{Post} explicitly, confirming this omission does not materially affect the result (Section~\ref{sec:res-baseline}).
 
\subsection{Item-stacked difference-in-differences}
\label{sec:item-stacked}
 
Fitting the ten deprivation items as ten independent regressions treats each item's error term as unrelated to the others, but the same child answers all ten deprivation questions, so their responses are correlated within household, a \citet{moulton1990} problem. As a check on this, the ten items are also reshaped into long format and fit as a single joint model with item and year fixed effects, clustering standard errors at household level to reflect the within-child correlation directly; full results are reported in Appendix~\ref{sec:appendix-stacked}.
 
\subsection{Double/debiased machine learning difference-in-differences}
\label{sec:dml}
 
Stage 3 asks whether the Section~\ref{sec:baseline-did} result depends on OLS's linear functional form and narrow six-variable control set, rather than on the treatment itself. A doubly robust, machine-learning-augmented DiD estimator relaxes both without reintroducing the bias a naive substitution would cause, for the reasons set out in Section~\ref{sec:lit-dml}.
 
The estimator used is \citet{zimmert2020}'s efficient-score, doubly robust DiD estimator, implemented via \texttt{causalweight::didDML()}, chosen over \citet{chang2020}'s alternative for the reasons given in Section~\ref{sec:lit-dml}. For each of the four treatment$\times$time cells, the estimator fits a cross-fitted outcome-regression model and a cell-propensity model, combined into a doubly robust, Neyman-orthogonal score (Section~\ref{sec:lit-dml}) that does not bias the ITT estimate at first order.\footnote{Cross-fitting uses $k=3$ folds with a 5\% propensity-score trim, both \texttt{causalweight::didDML()} defaults. Appendix Table~\ref{tab:dmlstability} tests sensitivity to the cross-fitting split directly, by refitting the preferred specification 20 times under different fold-assignment seeds.}
 
Two nuisance-model methods are compared: lasso, a linear method that (like OLS) cannot represent nonlinearities or interactions, and random forest, which can. Two covariate sets are also compared: the lean six CASE controls and the wide $\sim$35-variable set, feasible here only because of cross-fitting, since an earlier group-time ATT (\citeauthor{callawaysantanna2021}-style) attempt at the same wide set produced singular, rank-deficient models. Resulting ITT estimates for the official MDCH flag are reported under all combinations in Section~\ref{sec:res-dml}, plus a lasso/random-forest ensemble for the lean specification only, since the ensemble did not complete within a feasible runtime for the wide specification.

\subsection{Item-level and stacked machine-learning DiD}
\label{sec:item-dml}
 
The Section~\ref{sec:dml} estimator is repeated separately for each of the ten deprivation items, using the wide covariate set under both lasso and random forest nuisance models, with Benjamini--Hochberg correction applied across items. A fifth stage extends this to a single, jointly estimated model across all ten items using a shared complete-case sample and a multivariate outcome-regression learner, since no existing DML-DiD implementation accepts a vector-valued outcome directly; a hand-reconstructed single-item check confirmed the joint reimplementation is faithful. Both the per-item and the jointly estimated results are reported in Section~\ref{sec:results}.